\bookchapter{Bab 53: Penyimpanan Data dan Permintaan Rangkaian}

\chapterintro{Aplikasi yang hanya boleh 'dilihat' ialah aplikasi mati; aplikasi yang boleh menyimpan data dan menyambung rangkaian barulah hidup. Bab ini memperkenalkan 'set tiga' lapisan data Android: DataStore (menyimpan konfigurasi kecil), Room (menyimpan data berstruktur), Retrofit (berbual dengan pelayan), ditambah 'penyelaras utama' mereka — korutin. Selepas bab ini, anda boleh menulis aplikasi sebenar yang mengingati tabiat pengguna dan mengemas kini kandungan secara masa nyata.}

\booksection{Macam Mana Memilih Penyelesaian Penyimpanan: Lihat Jadual Ini Dahulu}

\begin{center}
\begin{tabularx}{\linewidth}{|X|X|X|X|}
\hline
\textbf{Penyelesaian} & \textbf{Sesuai menyimpan apa} & \textbf{Ciri} & \textbf{Diguna di mana dalam buku ini} \tabularnewline
\hline
DataStore & Konfigurasi berselerak (tema, suis) & Pasangan kunci-nilai; mudah dan selamat & Tetapan buku perakaunan Bab 54 \tabularnewline
Room & Data berstruktur (akaun, pengguna) & Pembungkus pangkalan data SQLite sebenar, boleh buat pertanyaan & Rekod akaun Bab 54 \tabularnewline
Retrofit & (Tidak disimpan, diambil melalui rangkaian) data pelayan & Piawaian de facto pustaka permintaan HTTP & Data cuaca Bab 55 \tabularnewline
\hline
\end{tabularx}
\end{center}

\begin{center}\small\color{gray}Pembahagian tugas set tiga data\end{center}

Hafal satu rumus: 'konfigurasi kecil masuk DataStore, data berkelompok masuk Room, data luar guna Retrofit'. SharedPreferences yang kerap muncul dalam tutorial lama sudah digantikan oleh DataStore, dan penulisan SQL SQLite secara manual juga sudah digantikan oleh Room — belajar terus yang baharu, jangan ikut jalan lama.

\booksection{DataStore: Menyimpan Sedikit Konfigurasi Kecil}

Mahu mengingati nilai berselerak seperti 'mata wang yang dipilih pengguna kali terakhir' atau 'mod gelap diaktifkan atau tidak', guna DataStore (versi Preferences):

\begin{lstlisting}[language=Kotlin]
val Context.dataStore by preferencesDataStore(name = "settings")

// Tulis: simpan satu nilai
suspend fun saveTheme(isDark: Boolean) {
    context.dataStore.edit { prefs ->
        prefs[THEME_KEY] = isDark
    }
}

// Baca: ambil nilainya (Flow memberitahu perubahan secara automatik)
val themeFlow: Flow<Boolean> = context.dataStore.data
    .map { it[THEME_KEY] ?: false }
\end{lstlisting}

Perhatikan yang dipulangkan ialah Flow — satu aliran data yang 'mengalir'. Sebaik sahaja nilai berubah, semua tempat yang memerhatinya akan menerima nilai baharu secara automatik. Konsep ini (reaktif) ialah batu asas Android moden, akan dipraktikkan dalam Bab 54.

\booksection{Room: Memberi Data 'Rumah Sebenar'}

Buku perakaunan ada ratusan hingga ribuan rekod, malah perlu pertanyaan dan statistik mengikut kategori dan bulan — ini memerlukan pangkalan data. Room ialah pembungkus rasmi SQLite, menjadikan 'membina jadual, membuat pertanyaan data' semula jadi seperti menulis Kotlin. Tiga benda:

\begin{lstlisting}[language=Kotlin]
// 1. Entity: struktur jadual (satu jadual = satu data class)
@Entity(tableName = "records")
data class Record(
    @PrimaryKey(autoGenerate = true) val id: Long = 0,
    val amount: Double,
    val category: String,
    val note: String,
    val date: String
)

// 2. Dao: antara muka operasi data (tambah, padam, ubah, cari)
@Dao
interface RecordDao {
    @Query("SELECT * FROM records ORDER BY date DESC")
    fun getAll(): Flow<List<Record>>

    @Insert
    suspend fun insert(record: Record)

    @Delete
    suspend fun delete(record: Record)

    @Query("SELECT SUM(amount) FROM records WHERE category = :c")
    suspend fun sumByCategory(c: String): Double
}

// 3. Database: pemasangan akhir
@Database(entities = [Record::class], version = 1)
abstract class AppDatabase : RoomDatabase() {
    abstract fun recordDao(): RecordDao
}
\end{lstlisting}

Templat 'set tiga' ini boleh dipindahkan terus ke mana-mana projek anda — tukar nama medan sahaja. Pertanyaan memulangkan Flow, antara muka segar semula mengikut data secara automatik; nikmatnya tidak terkata.

\booksection{Retrofit: Berbual dengan Pelayan}

Cuaca, berita, kadar pertukaran… data-data ini tiada dalam telefon, tetapi di pelayan orang lain. Retrofit menjadikan 'meminta data daripada pelayan' semudah memanggil fungsi. Tiga langkah:

\begin{lstlisting}[language=Kotlin]
// 1. Takrifkan antara muka: gambarkan permintaan HTTP sebagai fungsi
interface WeatherApi {
    @GET("v1/weather")
    suspend fun getWeather(@Query("city") city: String): WeatherResponse
}

// 2. Cipta instance
val api = Retrofit.Builder()
    .baseUrl("https://api.example.com/")
    .addConverterFactory(GsonConverterFactory.create())
    .build()
    .create(WeatherApi::class.java)

// 3. Panggil seperti memanggil fungsi (di dalam coroutine)
val result = api.getWeather("Beijing")
\end{lstlisting}

\textbf{Amaran Perangkap: Dua Hukum Besi Rangkaian}\par

(1) Permintaan rangkaian sama sekali tidak boleh dibuat pada utas utama — akan menyekat antara muka malah meranap (NetworkOnMainThreadException). Semua panggilan rangkaian mesti diletakkan dalam korutin (viewModelScope.launch). (2) HTTP berjelas dilarang secara lalai: mulai Android 9, HTTPS diwajibkan. Semasa menguji jika perlu guna alamat tempatan http://, perlu mengisytiharkan usesCleartextTraffic secara jelas dalam Manifest — tetapi produk sebenar mesti sepenuhnya HTTPS.

\booksection{Penghuraian JSON: 'Menterjemah' Data Rangkaian}

Data yang dipulangkan pelayan kebanyakannya JSON — rentetan teks seperti 'kamus tersarang'. Retrofit memerlukan seorang 'penterjemah' untuk menukarkannya kepada objek Kotlin. Dua penyelesaian arus perdana 2026:

\begin{center}
\begin{tabularx}{\linewidth}{|X|X|X|}
\hline
\textbf{Penyelesaian} & \textbf{Ciri} & \textbf{Senario disyorkan} \tabularnewline
\hline
Gson / Moshi & Penyelesaian refleksi veteran; kos sifar untuk konfigurasi & Projek kecil yang mahu bermula pantas, medan tidak banyak \tabularnewline
kotlinx.serialization & Penyelesaian semasa kompilasi rasmi; lebih pantas dan selamat & Pilihan utama projek baharu (arus perdana 2026) \tabularnewline
\hline
\end{tabularx}
\end{center}

\begin{center}\small\color{gray}Dua penyelesaian penghuraian JSON\end{center}

\begin{lstlisting}[language=Kotlin]
// Tambah anotasi @Serializable pada data class
@Serializable
data class WeatherResponse(
    val weather: List<Weather>,
    val main: Main,
    val name: String
)

// Retrofit memasang Converter berkenaan (contoh dengan sokongan rasmi Retrofit 2.11+):
// .addConverterFactory(json.asConverterFactory("application/json".toMediaType()))
\end{lstlisting}

Apa pun penyelesaiannya, hukum besinya ialah: lihat dahulu JSON sebenar yang dipulangkan, barulah takrifkan medan. Banyak pepijat 'penghuraian gagal' berpunca daripada nama medan tersilap atau tahap penyarangan tersalah jangka — cetak JSON mentah dan bandingkan; sepuluh saat dapat dikesan.

\booksection{Korutin: 'Bahasa Perantaraan' Keserentakan Android}

\booksection{Korutin: 'Bahasa Perantaraan' Keserentakan Android}

suspend dan launch yang berulang di atas datang daripada korutin Kotlin (Coroutine) — penyelesaian keserentakan rasmi Android. Anda boleh memahaminya secara kasar: memastikan operasi yang 'perlu menunggu sebentar' (rangkaian, pangkalan data) tidak menyekat antara muka.

\begin{lstlisting}[language=Kotlin]
// Memulakan coroutine dalam ViewModel
viewModelScope.launch {
    val records = dao.getAll().first()  // suspend: menunggu secara automatik tanpa menghalang
    // Kemas kini state paparan……
}

// Jembalitan kata kunci:
// suspend   —— "aku perlu tunggu, jangan halangi orang lain"
// launch    —— "Mula kerja!"
// viewModelScope —— "ruang kerja" yang mengikut halaman, berhenti sendiri apabila halaman dimusnahkan
\end{lstlisting}

Untuk yang sifar asas, cukup ingat cara guna dahulu; prinsip 'gantung dan sambung semula' di sebaliknya boleh dikaji lebih dalam kemudian. 90\% daripada apa yang anda tulis dalam pembangunan harian ialah templat ini.

\begin{tcbpractice}{Mari cuba: Buat Aplikasi Kecil 'Satu Ayat Sehari'}
Cari satu API petikan percuma (contohnya hitokoto dan seumpamanya), guna Retrofit untuk menarik 'satu ayat sehari', paparkan di tengah skrin, tambah butang 'Tukar ayat' (minta semula dalam korutin). Kemudian simpan ayat yang digemari ke dalam 'kegemaran' dengan Room. Latihan ini menggunakan DataStore/Room/Retrofit/korutin semuanya sekali habis — selepas menyiapkannya, Bab 7 dan 8 cuma makanan ringan.

\end{tcbpractice}

\begin{tcbsummary}{Ringkasan Bab}
\begin{itemize}
  \item Rumus: konfigurasi kecil masuk DataStore, data berkelompok masuk Room, data luar guna Retrofit
  \item DataStore memulangkan Flow: nilai berubah, pemerhati terus menerimanya secara automatik
  \item Set tiga Room: @Entity (jadual) + @Dao (operasi) + @Database (pemasangan akhir), templat boleh diguna semula terus
  \item Tiga langkah Retrofit: takrif antara muka cipta contoh panggil dalam korutin
  \item Dua hukum besi: rangkaian tidak masuk utas utama; produk sebenar mesti HTTPS
  \item Korutin = menunggu tanpa menyekat antara muka; viewModelScope ialah bengkel khusus halaman
\end{itemize}

\end{tcbsummary}
